\chapter{绪论}\label{chap:introduction}
\section{研究背景与意义}
博士论文需要明确研究对象、关键问题及其应用背景。绪论宜从问题出发组织叙述，并说明现有方法在目标条件下的不足。本文档中的文字用于说明模板的使用方式，不构成特定研究领域的论证。

\section{国内外研究现状}
引用文献时应确认所引文献能够支持相应论述。模板使用BibTeX管理文献，作者年份引用示例为\citet{knuth1984}，括号引用示例为\citep{lamport1994}。参考文献由发行版提供的GB/T 7714样式生成，不需要复制个人论文的文献库。

\subsection{问题分类与研究脉络}
按照研究问题梳理已有方法，比较其假设、输入信息、适用条件和局限。避免将文献综述写成与研究问题无关的方法清单。

\section{研究内容与章节安排}
第\ref{chap:foundations}章展示基本理论与符号；第\ref{chap:method}章给出方法、公式和矢量图示例；第\ref{chap:experiments}章展示实验表格与组合图；后续章节演示应用、扩展研究与总结的组织方式。实际论文应根据研究逻辑安排章节，而非机械沿用示例标题。
